\documentclass{article}
\usepackage{enumitem}
\usepackage{amsmath}
\begin{document}

\title{Chapter 16: Digestion and Absorption}
\date{}
\maketitle

\begin{enumerate}[label=Q\arabic*.]

\item The enzyme pepsin is active in the stomach because:
\\(A) Stomach pH is alkaline (pH 8--9)
\\(B) Stomach pH is highly acidic (pH 1.5--2) and pepsin has optimal activity at low pH
\\(C) HCl directly converts pepsin to pepsinogen
\\(D) Pepsin works at neutral pH only

\item Which enzyme is responsible for the conversion of inactive trypsinogen to active trypsin in the duodenum?
\\(A) Chymotrypsin
\\(B) Enterokinase (enteropeptidase)
\\(C) Pepsin
\\(D) Lipase

\item The brush border enzymes on microvilli of small intestine include all EXCEPT:
\\(A) Maltase
\\(B) Lactase
\\(C) Enterokinase
\\(D) Pepsinogen

\item Bile salts produced in the liver contribute to fat digestion by:
\\(A) Directly hydrolyzing triglycerides
\\(B) Emulsifying fats into smaller droplets, increasing surface area for lipase action
\\(C) Neutralizing stomach acid in the duodenum
\\(D) Activating pancreatic amylase

\item Assertion: Absorption of glucose in the small intestine is an active process.\\
Reason: Glucose is co-transported with Na$^+$ against its concentration gradient using energy from the Na$^+$/K$^+$ ATPase pump.
\\(A) Both true, Reason is correct explanation
\\(B) Both true, Reason is not correct explanation
\\(C) Assertion true, Reason false
\\(D) Assertion false, Reason true

\item The villi and microvilli in the small intestine increase absorption efficiency by:
\\(A) Secreting digestive enzymes
\\(B) Increasing the surface area for absorption enormously
\\(C) Producing mucus for lubrication
\\(D) Peristaltic movement of food

\item Intrinsic factor secreted by gastric parietal cells is essential for absorption of:
\\(A) Vitamin C
\\(B) Vitamin B$_{12}$ (cobalamin)
\\(C) Iron
\\(D) Calcium

\item Fat-soluble vitamins (A, D, E, K) are absorbed via:
\\(A) Portal blood directly to liver
\\(B) Lacteal lymph vessels in villi (as chylomicrons) to lymph system
\\(C) Active transport through enterocytes
\\(D) Pinocytosis only

\item The cephalic phase of gastric secretion is triggered by:
\\(A) Food entering the stomach
\\(B) Sight, smell, and thought of food via the vagus nerve
\\(C) Secretin released by duodenal cells
\\(D) Cholecystokinin (CCK) released in intestine

\item Peristalsis in the digestive tract is controlled primarily by:
\\(A) Voluntary control by the somatic nervous system
\\(B) The enteric nervous system (myenteric plexus) within the gut wall
\\(C) Hormones from pancreas only
\\(D) Circular muscle contractions only

\item Which hormone stimulates the pancreas to secrete enzyme-rich pancreatic juice?
\\(A) Secretin
\\(B) Cholecystokinin (CCK)
\\(C) Gastrin
\\(D) Motilin

\item Deficiency of lactase enzyme leads to lactose intolerance because:
\\(A) Lactose is toxic to gut cells
\\(B) Undigested lactose in the colon is fermented by bacteria, producing gas and osmotic diarrhea
\\(C) Lactose inhibits glucose absorption
\\(D) Lactose damages intestinal villi

\item In the large intestine, which of the following primarily occurs?
\\(A) Protein digestion
\\(B) Fat absorption
\\(C) Water and electrolyte absorption, and bacterial synthesis of vitamin K
\\(D) Glucose absorption

\item Match the following digestive enzymes with their substrates:\\
1. Amylase $\rightarrow$ (a) Starch\\
2. Trypsin $\rightarrow$ (b) Proteins\\
3. Lipase $\rightarrow$ (c) Triglycerides\\
4. Nuclease $\rightarrow$ (d) Nucleic acids
\\(A) 1-a, 2-b, 3-c, 4-d
\\(B) 1-b, 2-a, 3-d, 4-c
\\(C) 1-c, 2-d, 3-a, 4-b
\\(D) 1-d, 2-c, 3-b, 4-a

\item The pyloric sphincter regulates:
\\(A) Entry of food from esophagus to stomach
\\(B) Passage of partially digested food (chyme) from stomach to duodenum
\\(C) Bile entry from bile duct to duodenum
\\(D) Movement of food in the colon

\end{enumerate}
\end{document}
